\section{Introduction}

Digit-distribution tests ask whether reported election counts have unusual
leading, second, or final digits. They are attractive because they are simple:
a count table can be reduced to a reference distribution, a test statistic, and
a ranked list of contests or jurisdictions that deserve a closer look. In the
W-series, that is all they are allowed to do. They can prioritize
investigation. They cannot prove fraud, certify correctness, or replace a
risk-limiting audit.

This paper also leads with a second boundary specific to this method family.
The election-forensics literature is skeptical of Benford's law for vote
counts. Benford behavior depends on measurements that span several orders of
magnitude and are not tightly bounded by administrative design. Precinct vote
counts often do the opposite: jurisdictions build precincts to be similar in
size, contests are bounded by registration and turnout, and many tables cover a
narrow range of magnitudes. First-digit tests are therefore fragile, and
second-digit tests are especially easy to misread. Deviations from Benford-style
references are expected in many innocent election tables.

The purpose of W.04 is to define a cautious method wrapper for these weak
screens. Following the W.01 framework contract \citep{w01-forensic-analytics-2026},
it specifies the digit-test families, the RCOUNT inputs they consume, the
transcript fields they must emit, and the fixture roles that keep both planted
anomalies and innocent failures visible. It is a method and framework paper. No
forensics crate, function, fixture, or test is claimed to exist here.
